\section{Open problems}\label{sec:open}

We list what the paper leaves open, in the order of the sections.
\begin{enumerate}[(1)]
\item (The start of the big-jump zone.) When the plus scale is the smaller one, the window
$KB_N\le k-pN\le K\lambda_NB_N$ lies between Theorems~\ref{thm:bulk} and~\ref{thm:bigjump}.
Theorem~\ref{thm:window} locates the start of the zone up to a factor $1\pm\delta$ for fresh
starts and describes $P_N$ below it on the logarithmic scale. A relative-error statement on
the whole window, the strip $(1\pm\delta)\xi_N$ around the start, and the sign of the first
differences on the window (Remark~\ref{rem:mode-open}) are open. Since (EQ) excludes this
case, the exact mode count is open whenever the two scales are not comparable.
\item (Every $N$.) Theorem~\ref{thm:mode-count} holds for all large $N$, and its threshold
depends on the laws. The count for the law of Example~\ref{ex:runs2} is $17$ or $19$ at
$N\le2000$ and $15$ from $N=2500$ on. We know no bound on the number of modes that holds for
every $N$, even for the fair start of a log-convex law.
\item (Total-progeny laws.) Corollary~\ref{cor:exposure-mode}(c) needs each end profile to
have one local maximum, which holds for strictly log-convex cycle laws. For the cycle law of
the slotted queue with zeta service we have checked strict log-convexity numerically up to
$n=2^{14}$. A proof, or a direct proof that $u_C$ is decreasing, would make the 3 modes
unconditional.
\item (The line $\alpha=\varphi$.) On the line $\rho_\sgn=1/\alpha_*$ of
Theorem~\ref{thm:mode}(c) the comparison of the endpoint with the bulk depends on slowly
varying factors. For pure power laws and the stationary or aged start the constant is
explicit and below $1$. For the cycle laws of Section~\ref{sec:exposure} it is not computed.
\item (The boundary indices.) At $\alpha=1$ the runs may or may not have a finite mean, and the
scale $B_N$ is $N$ up to a slowly varying factor. The local tools of Berger~\cite{berger2019ptrf}
for the Cauchy domain should give the bulk and the big-jump zone, with the centering
$pN$ replaced by a truncated mean. At $\alpha=2$ the law has a Gaussian bulk with a slowly
varying variance and a big-jump zone that starts at a logarithmic multiple of the scale. We
have not worked out either line. The case $0<\alpha<1$ of infinite means, where the bulk is
the Lamperti law and the scale is $N$, is the subject of~\cite{paperG}.
\end{enumerate}
